\documentclass[twocolumn, 12pt]{article}

\usepackage[utf8]{inputenc}
\usepackage[english, spanish]{babel}
\usepackage{fullpage}
\usepackage{graphicx}
\usepackage{amsmath}
\usepackage{enumitem}
\usepackage{chngcntr}
\usepackage{setspace}
\usepackage{url}
\usepackage{csquotes}
\usepackage{float}
\usepackage{verbatim}
\usepackage{tabularx}
\usepackage{amsmath}
\usepackage{caption}
\usepackage{bm}
\usepackage{colortbl}
\usepackage{xcolor}

\usepackage{multirow}

% \usepackage{hyperref}

\definecolor{LigthGray}{rgb}{0.7098, 0.7294, 0.7215}
\definecolor{LigthGrayPlus}{rgb}{0.6862, 0.7254, 0.7294}
\definecolor{White}{rgb}{1, 1, 1}
\definecolor{Red}{rgb}{1, 0, 0}

\counterwithin{figure}{section}
\renewcommand{\thesection}{\arabic{section}}
\renewcommand{\thesubsection}{\thesection.\arabic{subsection}}
\renewcommand{\baselinestretch}{1.5}

\usepackage[style=apa, maxnames=6, minnames=3, backend=biber]{biblatex}
\DefineBibliographyStrings{english}{%chktex-file 1 chktex-file 6
    andothers = {\em et\addabbrvspace al\adddot}
}
\addbibresource{./Bibliography/bibliography.bib}

\usepackage{array}

\setlength{\parskip}{0pt}

\raggedbottom{}

\newcommand{\bolditalic}[1]{\textbf{\textit{#1}}}

\begin{document}

\begin{titlepage}
    \centering
    \includegraphics[width=0.3\textwidth]{Images/logo_utb.png}\par\vspace{1cm}
    {\scshape\LARGE Universidad Tecnológica de Bolívar \par}
    \vspace{1cm}

    {\scshape\Large FÍSICA CALOR Y ONDAS \par}
    \vspace{.2cm}

    % chktex-file 8
    {\scshape\Large Grupo 1 \par}
    \vspace{1cm}
    % chktex-file 8
    \slshape {\Large \bfseries{}Informe de Laboratorio No. II\\}
    \slshape {\small \bfseries{} ONDAS ESTACIONARIAS EN UNA CUERDA.~RESONANCIA:\@{}~MOVIMIENTO ARMÓNICO SIMPLE.}
    \vspace{2cm}

    \slshape {\itshape{} Mauro González, T00067622 \\}
    \slshape {\itshape{} German De Armas Castaño, T00068765 \\}
    \slshape {\itshape{} Angel Vega Rodriguez, T00068186 \\}
    \slshape {\itshape{} Juan Jose Osorio Ariza, T00067316 \\}
    \slshape {\itshape{} Jorge Alberto Rueda Salgado, T00068722 \\}
    \vfill
    Revisado Por \\
    Duban Andres Paternina Verona\\
    {\large \today\par}
\end{titlepage}

% chktex-file 44

% ! ----------------------------------------------------------------------|>
\section{Introducción}

Las ondas estacionarias son un tipo especial de ondas que se forman cuando dos ondas 
viajeras de igual frecuencia y amplitud se superponen y se mueven en direcciones opuestas. 
En una onda estacionaria, los puntos individuales que la componen experimentan oscilaciones 
alrededor de sus posiciones de equilibrio a medida que pasa el tiempo. Sin embargo, lo que 
distingue a una onda estacionaria es que su patrón general no se desplaza a lo largo del medio, 
de ahí su denominación.

En esta experiencia de laboratorio, nos centraremos en la exploración y comprobación de 
cada uno de los conceptos relacionados con la propagación de las ondas estacionarias en 
una cuerda. Para ello, se utilizará una cuerda flexible y una fuente de 
vibración controlable que nos permitirán observar y analizar cómo cambian 
los patrones de ondas estacionarias.


% ! ----------------------------------------------------------------------|>
\section{Objetivos}

% + ----------------------------------------|>
\subsection{Objetivo general}

\begin{itemize}[label=$\triangleright$]
    \item Explicar las propiedades de las ondas estacionarias en un contexto 
        experimental específico, a partir del análisis de factores como la 
        frecuencia, la longitud de onda y la tensión en el medio.    
\end{itemize}

% + ----------------------------------------|>
\subsection{Objetivos específicos}

\begin{itemize}[label=$\triangleright$]
    \item Determinar las propiedades de una onda estacionaria tales como la 
        longitud de onda, la amplitud, la frecuencia y la velocidad de una onda 
        estacionaria en un medio específico.

    \item Identificar los puntos de formación de nodos (puntos de amplitud mínima) 
        y antinodos (puntos de amplitud máxima).

    \item Estudiar la relación entre la frecuencia y la longitud de onda y el como 
        afectan la formación y el patrón de las ondas estacionarias en una cuerda
\end{itemize}

% ! ----------------------------------------------------------------------|>
\section{Marco Teórico}

% * -------------------------------------------------------|>
\subsection{Onda mecánica}

Las ondas mecánicas avanzan a través de un medio elástico,
cuyas partículas oscilan en torno a un punto fijo. El medio
en cuestión puede ser gaseoso, líquido o sólido.

Para que exista una onda mecánica es necesario que haya una
fuente que genere la perturbación y un medio por el cual
dicha perturbación pueda propagarse. Además se necesita un
medio físico que permita a los elementos influirse entre
sí.~(\cite{J_Merino_2022})

% * -------------------------------------------------------|>
\subsection{Expresión para calcular la velocidad de una onda en una cuerda}

La velocidad de un pulso u onda en una cuerda bajo tensión
se puede calcular con la ecuación:

\begin{equation}
    v = \sqrt{\frac{F_T}{\mu}}
    \label{eq:velocidad-onda-cuerda}
\end{equation}

donde $F_T$ es la tensión de la cuerda y $\mu$ es la masa
por longitud de la cuerda.~(\cite{Moebs_2021})

% * -------------------------------------------------------|>
\subsection{Formulas para ondas estacionarias en una cuerda}

Para una cuerda fija en ambos extremos que forma una onda
estacionaria, los modos normales de vibración se
caracterizan por diferentes longitudes de onda y
frecuencias.

\begin{itemize}[label=$\triangleright$]
    \item Longitud de onda de modos normales: La longitud de onda
          para el n-ésimo modo normal de vibración en una cuerda fija
          en ambos extremos está dada por:

          \begin{equation}
              \lambda_{n} = \frac{2 \cdot L}{n}
              \label{eq:n_lambda}
          \end{equation}

    \item Frecuencia de modos normales en función de la rapidez de la
          onda y longitud de la cuerda: La frecuencia $f_n$ del
          n-ésimo modo normal está relacionada con la rapidez de la
          onda $v$ y la longitud de onda ($\lambda_{n}$).

          \begin{equation*}
              f_n = \frac{v}{\lambda_{n}} = \frac{n v}{2L}
          \end{equation*}

    \item Frecuencia de modos normales en función de la tensión de la
          cuerda y la densidad lineal de masa: La frecuencia del
          n-ésimo modo normal también está relacionada con la tensión
          en la cuerda Y la densidad lineal de masa ($\mu$).

          \begin{equation}
              f = \frac{n}{2 L} \sqrt{\frac{F_T}{\mu}}
              \label{eq:n_frecuencia}
          \end{equation}

    \item Frecuencia fundamental: La frecuencia fundamental ($f_{1}$)
          s la frecuencia del primer modo normal (modo fundamental),
          que es cuando la cuerda vibra en una única mitad de ciclo.
          Para el modo fundamental $(n=1)$:

          \begin{equation*}
              f_{1} = \frac{v}{2L} = \frac{v}{2L}
          \end{equation*}

          Esta es la frecuencia más baja en la cual la cuerda puede
          vibrar.
\end{itemize}

% * -------------------------------------------------------|>
\subsection{Calculo de errores}

{\large
    \begin{equation}
        \begin{gathered}
            E_{absoluto} = \left\lvert V_E - V_A\right\rvert \\
            E_{relativo} = \frac{E_A}{V_E} \\
        \end{gathered}
        \label{eq:calculo-errores}
    \end{equation}
}

% ! ----------------------------------------------------------------------|>
\section{Montaje Experimental}




% ! ----------------------------------------------------------------------|>
\section{Datos Experimentales}

\begin{table}[H]
    \begin{center}
        \begin{tabularx}{0.9\linewidth}{|>{\centering\arraybackslash}X|>{\centering\arraybackslash}X|}
            \hline
            \multicolumn{2}{|c|}{\textbf{Constantes}}   \\\hline
            $M_{cuerda}$ & $1 \times 10^{-3} Kg$        \\\hline
            $L_{cuerda}$ & $1,425 m$                    \\\hline
            $\mu$        & $7,0175 \times 10^{-4} Kg/m$ \\\hline
            $L$          & $8,6 \times 10^{-1} m$       \\\hline
        \end{tabularx}
    \end{center}
\end{table}

\begin{table}[H]
    \begin{center}
        \begin{tabularx}{0.9\linewidth}{|>{\centering\arraybackslash}X|>{\centering\arraybackslash}X|}
            \hline
            \multicolumn{2}{|c|}{$M_{1} = 0,0349 Kg$}       \\\hline
            No. & Frecuencia experimental \bolditalic{(Hz)} \\\hline
            1   & $15,43$                                   \\\hline
            2   & $30,39$                                   \\\hline
            3   & $45,73$                                   \\\hline
            4   & $60,00$                                   \\\hline
        \end{tabularx}
    \end{center}
\end{table}

\vspace{-.5cm}

\begin{table}[H]
    \begin{center}
        \begin{tabularx}{0.9\linewidth}{|>{\centering\arraybackslash}X|>{\centering\arraybackslash}X|}
            \hline
            \multicolumn{2}{|c|}{$M_{2} = 0,0399 Kg$}       \\\hline
            No. & Frecuencia experimental \bolditalic{(Hz)} \\\hline
            1   & $12,98$                                   \\\hline
            2   & $26,00$                                   \\\hline
            3   & $39,24$                                   \\\hline
            4   & $52,20$                                   \\\hline
        \end{tabularx}
    \end{center}
\end{table}

\vspace{-.5cm}

\begin{table}[H]
    \begin{center}
        \begin{tabularx}{0.9\linewidth}{|>{\centering\arraybackslash}X|>{\centering\arraybackslash}X|}
            \hline
            \multicolumn{2}{|c|}{$M_{3} = 0,0449 Kg$}       \\\hline
            No. & Frecuencia experimental \bolditalic{(Hz)} \\\hline
            1   & $13,56$                                   \\\hline
            2   & $28,50$                                   \\\hline
            3   & $43,04$                                   \\\hline
            4   & $58,38$                                   \\\hline
        \end{tabularx}
    \end{center}
\end{table}

% ! ----------------------------------------------------------------------|>
\section{Análisis de datos}

% + ----------------------------------------|>
\subsection{Análisis}

% * -----------------------------|>
\subsubsection{Calculo de errores}

Usando la formula~(\ref{eq:n_frecuencia})
y~(\ref{eq:calculo-errores}),

\begin{table}[H]
    \begin{tabularx}{.9\linewidth}{|>{\centering\arraybackslash}X|>{\centering\arraybackslash}X|>{\centering\arraybackslash}X|}
        \hline
        \multicolumn{3}{|c|}{$M_{1}$}                                        \\\hline
        No. & Frecuencia Teórica \bolditalic{(Hz)} & Error \bolditalic{(\%)} \\\hline
        1   & $12,8418$                            & $20,1544$               \\\hline
        2   & $25,6836$                            & $18,3245$               \\\hline
        3   & $38,5254$                            & $18,7008$               \\\hline
        4   & $51,3672$                            & $16,8060$               \\\hline

    \end{tabularx}
\end{table}

\vspace{-.5cm}

\begin{table}[H]
    \begin{tabularx}{.9\linewidth}{|>{\centering\arraybackslash}X|>{\centering\arraybackslash}X|>{\centering\arraybackslash}X|}
        \hline
        \multicolumn{3}{|c|}{$M_{2}$}                                        \\\hline
        No. & Frecuencia Teórica \bolditalic{(Hz)} & Error \bolditalic{(\%)} \\\hline
        1   & $13,7309$                            & $5,4689$                \\\hline
        2   & $27,4619$                            & $5,3232$                \\\hline
        3   & $41,1928$                            & $4,7406$                \\\hline
        4   & $54,9237$                            & $4,9591$                \\\hline
    \end{tabularx}
\end{table}

\vspace{-.5cm}

\begin{table}[H]
    \begin{tabularx}{.9\linewidth}{|>{\centering\arraybackslash}X|>{\centering\arraybackslash}X|>{\centering\arraybackslash}X|}
        \hline
        \multicolumn{3}{|c|}{$M_{3}$}                                        \\\hline
        No. & Frecuencia Teórica \bolditalic{(Hz)} & Error \bolditalic{(\%)} \\\hline
        1   & $14,5659$                            & $6,9057$                \\\hline
        2   & $29,1318$                            & $2,1686$                \\\hline
        3   & $43,6976$                            & $1,5050$                \\\hline
        4   & $58,2635$                            & $0,1999$                \\\hline

    \end{tabularx}
\end{table}

% * -----------------------------|>
\subsubsection{Rapidez de propagacion en funcion de la frecuencia}

Usando la formula~(\ref{eq:n_lambda}) y

\begin{equation*}
    \lambda = \frac{v}{f}
\end{equation*}

\begin{table}[H]
    \begin{tabularx}{.9\linewidth}{|>{\centering\arraybackslash}X|}
        \hline
        $\lambda_{n}$ \bolditalic{(m)} \\\hline
        $1: 1,7200$                    \\\hline
        $2: 0,8600$                    \\\hline
        $3: 0,5733$                    \\\hline
        $4: 0,4300$                    \\\hline
    \end{tabularx}
\end{table}

\begin{table}[H]
    \begin{tabularx}{0.9\linewidth}{|>{\centering\arraybackslash}X|>{\centering\arraybackslash}X|>{\centering\arraybackslash}X|}
        \multicolumn{3}{c}{\large $M_{1}$}                                                                         \\\hline
        Velocidad Experimental \bolditalic{(m/s)} & Velocidad Teórica \bolditalic{(m/s)} & Error \bolditalic{(\%)} \\\hline
        $26,5396$                                 & $22,0879$                            & $20,1544$               \\\hline
        $26,1354$                                 & \dots                                & $18,3245$               \\\hline
        $26,2185$                                 & \dots                                & $18,7008$               \\\hline
        $25,8000$                                 & \dots                                & $16,8060$               \\\hline
    \end{tabularx}
\end{table}

\vspace{-.5cm}

\begin{table}[H]
    \begin{tabularx}{0.9\linewidth}{|>{\centering\arraybackslash}X|>{\centering\arraybackslash}X|>{\centering\arraybackslash}X|}
        \multicolumn{3}{c}{\large $M_{2}$}                                                                         \\\hline
        Velocidad Experimental \bolditalic{(m/s)} & Velocidad Teórica \bolditalic{(m/s)} & Error \bolditalic{(\%)} \\\hline
        $22,3256$                                 & $23,6172$                            & $5,4689$                \\\hline
        $22,3600$                                 & \dots                                & $5,3232$                \\\hline
        $22,4976$                                 & \dots                                & $4,7406$                \\\hline
        $22,4460$                                 & \dots                                & $4,9591$                \\\hline
    \end{tabularx}
\end{table}

\vspace{-.5cm}

\begin{table}[H]
    \begin{tabularx}{0.9\linewidth}{|>{\centering\arraybackslash}X|>{\centering\arraybackslash}X|>{\centering\arraybackslash}X|}
        \multicolumn{3}{c}{\large $M_{3}$}                                                                         \\\hline
        Velocidad Experimental \bolditalic{(m/s)} & Velocidad Teórica \bolditalic{(m/s)} & Error \bolditalic{(\%)} \\\hline
        $23,3232$                                 & $25,0533$                            & $6,9057$                \\\hline
        $24,5100$                                 & \dots                                & $2,1686$                \\\hline
        $24,6763$                                 & \dots                                & $1,5050$                \\\hline
        $25,1034$                                 & \dots                                & $0,1999$                \\\hline
    \end{tabularx}
\end{table}

% * -----------------------------|>
\subsubsection{}

Al obtener una frecuencia de resonancia mayor, dada la
formula

\begin{equation*}
    \lambda = \frac{v}{f}
\end{equation*}

Esto quiere decir que $f$ y $\lambda$ son inversamente
proporcionales.

% * -----------------------------|>
\subsubsection{}

No depende la velocidad de propagación de una onda en una
cuerda de la frecuencia de la misma. Como
establece~(\ref{eq:velocidad-onda-cuerda}), la velocidad
dependerá solamente de la tension de la cuerda, ademas de
su densidad lineal.

% * -----------------------------|>
\subsubsection{}

Al aumentar la tension, también aumentara la velocidad de
propagación, asi como lo establece la
ecuación~(\ref{eq:velocidad-onda-cuerda})

% * -----------------------------|>
\subsubsection{}

Al aumentar la densidad lineal de la cuerda, la velocidad
de propagación disminuirá, debido a que son inversamente
proporcionales como lo
establece~(\ref{eq:velocidad-onda-cuerda})

% * -----------------------------|>
\subsubsection{}

Si se mantiene la misma tension en la cuerda, pero se
disminuye la distancia entre los dos extremos, la
frecuencia fundamental aumentará.

Precisamente en la ecuación~(\ref{eq:n_frecuencia}), se
establece que la longitud y la frecuencia resultante son
inversamente proporcionales.

% ! ----------------------------------------------------------------------|>
\section{Conclusiones}

\printbibliography

\end{document}